%%%%%%%%%%%%%%%%%%%%%%%%%%%%%%%%%%%%%%%%%%%%%%%%%%%%%%%%%%%%%%%%%%%%%%
% Lớp 11 | Xác suất | Biến cố hợp và quy tắc cộng xác suất
% Dạng 1D9?2-5: Tính xác suất biến cố hợp bằng sơ đồ hình cây
% Số câu: 15  (sinh tự động bởi tach_cau.py; từ nay thầy sửa trực tiếp file này)
%%%%%%%%%%%%%%%%%%%%%%%%%%%%%%%%%%%%%%%%%%%%%%%%%%%%%%%%%%%%%%%%%%%%%%

%%%%%%%%%%%%%%%% TRẮC NGHIỆM 4 PHƯƠNG ÁN %%%%%%%%%%%%%%%%

% nguồn: Ok_TracNghiemLop11_Lan1
\begin{ex}%[1D9N2-5]%[ma:7ed98e79]
	Một trường có $50$ em học sinh giỏi trong đó có $5$ cặp anh em sinh đôi. Cần chọn ra $3$ học sinh trong số $50$  học sinh để tham gia trại hè. Tính xác suất trong $3$ em ấy không có cặp anh em sinh đôi nào.
	 \choice
		 {$\dfrac{102}{245}$}
		 {$\dfrac{48}{49}$}
		{$\dfrac{97}{98}$} 
		 {\True $\dfrac{242}{245}$}
	\loigiai{ 
		Số phần tử của không gian mẫu $n(\Omega)=\mathrm{C}_{50}^3$.\\ 
		Xét biến cố $A\colon$ \lq\lq Trong $3$ em ấy không có cặp anh em sinh  đôi\rq\rq, biến cố đối $\overline{A}\colon$ \lq\lq Trong $3$ em ấy có $1$ cặp anh em sinh đôi\rq\rq.\\
		Số phần tử của biến  cố $\overline{A}$ là $n(\overline{A})=\mathrm{C}_5^1\cdot \mathrm{C}_{48}^1$.\\
		Xác suất của biến cố  $\overline{A}$ là $\mathrm{P}(\overline{A})=\dfrac{\mathrm{C}_5^1\cdot \mathrm{C}_{48}^1}{\mathrm{C}_{50}^3}=\dfrac{3}{245}$.\\
	Vậy $\mathrm{P}(A)=1-\mathrm{P}(\overline{A})=\dfrac{242}{245}$.} 
\end{ex}

% nguồn: Ok_TracNghiemLop11_Lan1
\begin{ex}%[1D9N2-5]%[ma:382bd000]
	Cho $A$ và $B$ là hai biến  cố độc lập. Biết $\mathrm{P}(A)=0{,}4$ và $\mathrm{P}(B)=0{,}5$. Xác suất của biến cố $A\cup B$ là
	\choice 
		 {$0{,}9$}
		{\True  $0{,}7$}
		{$0{,}5$} 
		 {$0{,}2$}
	 \loigiai{
		Vì $A$ và $B$ là hai biến cố độc lập nên $\mathrm{P}(AB)=\mathrm{P}(A)\cdot  \mathrm{P}(B)=0{,}4\cdot0{,}5=0{,}2.$\\
	Do đó $\mathrm{P}(A\cup  B)=\mathrm{P}(A)+\mathrm{P}(B)-\mathrm{P}(AB)=0{,}4+0{,}5-0{,}2=0,7$.}
\end{ex}

% nguồn: Ok_TracNghiemLop11_Lan1
\begin{ex}%[1D9H2-5]%[ma:7ea68f7e]
	Cho $A$, $B$ là hai biến cố độc lập. Biết $\mathrm{P}(A)=0{,}5$; $\mathrm{P}(AB)=0{,}2$. Tính $\mathrm{P}(A\cup  B)$.
	 \choice
		 {$0{,}4$}
		 {$0{,}9$}
		 {$0{,}6$}
		{\True $0{,}7$} 
	\loigiai{ 
		Vì $A$, $B$ là hai  biến cố độc lập nên $\mathrm{P}(AB)=\mathrm{P}(A)\cdot \mathrm{P}(B)\Rightarrow \mathrm{P}(B)=\dfrac{\mathrm{P}(AB)}{\mathrm{P}(A)}=0{,}4$.\\
	 Do đó $\mathrm{P}(A\cup B)=\mathrm{P}(A)+\mathrm{P}(B)-\mathrm{P}(AB)=0{,}7$.}
\end{ex}

% nguồn: Ok_TracNghiemLop11_Lan1
\begin{ex}%[1D9V2-5]%[ma:d46b4c50]
	Một lớp học có $30$ học sinh, trong đó có $18$ học sinh thích đá bóng, $16$ học sinh thích cầu lông và $7$ học sinh thích chơi cả hai môn thể thao.  Chọn ngẫu nhiên một học sinh. Xác suất bạn học sinh được chọn không thích chơi môn thể thao nào là
	\choice 
		 {\True $\dfrac{1}{10}$}
		{$\dfrac{3}{5}$} 
		 {$\dfrac{7}{30}$}
		 {$\dfrac{1}{15}$}
	\loigiai{ 
		Ta có  $n(\Omega)=\mathrm{C}_{30}^1$.\\
		Ta gọi các  biến cố
		\begin{itemize}
			\item $A\colon$\lq\lq Bạn  học sinh được chọn thích chơi đá bóng\rq\rq.\\
			Suy  ra $n(A)=\mathrm{C}_{18}^1$ hay $\mathrm{P}(A)=\dfrac{18}{30}=\dfrac{3}{5}$;
			\item $\overline{A}\colon$\lq\lq Bạn học sinh  được chọn không thích chơi đá bóng\rq\rq.
			\item $B\colon$\lq\lq Bạn học sinh được chọn thích chơi cầu lông\rq\rq.\\ 
			 Suy ra $n(B)=\mathrm{C}_{16}^1$ hay $\mathrm{P}(B)=\dfrac{16}{30}=\dfrac{8}{15}$.
			\item $\overline{B}\colon$\lq\lq Bạn học sinh được chọn không  thích chơi cầu lông\rq\rq.
			\item $A\cap  B\colon$\lq\lq Bạn học sinh được chọn thích chơi cả hai môn thể thao\rq\rq.\\
			Suy ra $n(A\cap  B)=\mathrm{C}_7^1$ hay $\mathrm{P}(A\cap B)=\dfrac{7}{30}$.
			\item $\overline{A}\cap\overline{B}=\overline{A\cup B}\colon$\lq\lq Bạn học sinh  được chọn không thích chơi cả $2$ môn thể thao\rq\rq.
		\end{itemize}
		Ta có $\mathrm{P}(A\cup B)=\mathrm{P}(A)+\mathrm{P}(B)-\mathrm{P}(A\cap B)=\dfrac{9}{10}$.\\ 
	Suy ra $\mathrm{P}(\overline{A\cup B})=1-\mathrm{P}(A\cup  B)=\dfrac{1}{10}$.}
\end{ex}


%%%%%%%%%%%%%%%% ĐÚNG/SAI %%%%%%%%%%%%%%%%

% nguồn: Ok_dungsai_lop11_Lan1
\begin{ex}%[1D9H2-5]%[ma:fefcd779]
 Một hộp đựng $30$ tấm thẻ có đánh số từ $1$ đến $30$, hai tấm thẻ khác nhau  đánh hai số khác nhau. Lấy ngẫu nhiên một tấm thẻ từ hộp.
 \choiceTF 
  {\True Xác suất thẻ đánh số  chia hết cho $3$ bằng $\dfrac{1}{3}$}
  {Xác suất thẻ đánh số chia hết cho $4$ bằng $\dfrac{11}{30}$} 
  {\True Xác suất thẻ đánh số chia hết cho $3$ và chia  hết cho $4$ bằng $\dfrac{1}{15}$}
   {\True Xác suất thẻ đánh số chia hết cho $3$ hoặc $4$ bằng $\dfrac{1}{2}$}
  \loigiai{
  \begin{itemchoice}
    \itemch Gọi $A$ là biến cố \lq\lq Lấy được thẻ đánh số chia hết cho $3$\rq\rq.\\
   Suy ra $n(A)=10$ và  $\mathrm{P}(A)=\dfrac{10}{30}=\dfrac{1}{3}$.
   \itemch Gọi $B$ là  biến cố \lq\lq Lấy được thẻ đánh số chia hết cho $4$\rq\rq.\\
   Suy ra $n(B)=7$ và  $\mathrm{P}(B)=\dfrac{7}{30}$.
   \itemch Ta có $A B$ là biến cố \lq\lq Lấy được thẻ đánh số chia hết  cho $3$ và chia hết cho $4$\rq\rq.\\
    Suy ra $A B=\{12;24\}$ hay $n(A B)=2$ và $\mathrm{P}(A B)=\dfrac{2}{30}=\dfrac{1}{15}$.
   \itemch Xác suất để lấy  được thẻ đánh số chia hết cho $3$ hoặc $4$ là\\
   $\mathrm{P}(A\cup B)=\mathrm{P}(A)+\mathrm{P}(B)-\mathrm{P}(AB)=\dfrac{1}{3}+\dfrac{7}{30}-\dfrac{1}{15}=\dfrac{1}{2}$. 
  \end{itemchoice}
  }
\end{ex}

% nguồn: Ok_dungsai_lop11_Lan1
\begin{ex}%[1D9H2-5]%[ma:889edb56]
 Một lớp học có $40$ học sinh, trong đó có $18$ học sinh tham gia môn bóng đá và $10$ học sinh tham gia môn bóng chuyền, trong đó có $6$ học sinh tham gia cả hai môn bóng đá và bóng chuyền. Thầy giáo chọn ngẫu nhiên một học sinh từ lớp học để làm nhiệm vụ đặc biệt, gọi $A$ là biến cố \lq\lq Chọn được một học sinh tham gia môn bóng đá\rq\rq, $B$ là biến cố \lq\lq Chọn được một học sinh tham gia môn bóng chuyền\rq\rq. Khi  đó
  \choiceTF
  {\True  $\mathrm{P}(A)=\dfrac{9}{20}$}
  {\True  $\mathrm{P}(B)=\dfrac{1}{4}$}
   {$\mathrm{P}(AB)=\dfrac{7}{20}$}
  {Xác suất  để học sinh được chọn có tham gia ít nhất một trong hai môn thể thao bằng $\dfrac{13}{20}$}
 \loigiai{ 
  Số  phần tử của không gian mẫu $n(\Omega)=40$.
  \begin{itemchoice}
   \itemch Có $n(A)=18$ suy ra  $\mathrm{P}(A)=\dfrac{18}{40}=\dfrac{9}{20}$.
   \itemch  Có $n(B)=10$ suy ra $\mathrm{P}(B)=\dfrac{10}{40}=\dfrac{1}{4}$.
   \itemch $AB$ là biến cố chọn được học sinh tham gia cả hai môn.  Ta có $n(AB)=6$ suy ra $\mathrm{P}(A B)=\dfrac{6}{40}=\dfrac{3}{20}$.
   \itemch Xác suất để  chọn được một học sinh tham gia ít nhất một trong hai môn bóng đá, bóng chuyền là $\mathrm{P}(A\cup B)=\mathrm{P}(A)+\mathrm{P}(B)-\mathrm{P}(A B)=\dfrac{9}{20}+\dfrac{1}{4}-\dfrac{3}{20}=\dfrac{11}{20}$.
  \end{itemchoice}
 } 
\end{ex}

% nguồn: Ok_dungsai_lop11_Lan1
\begin{ex}%[1D9H2-5]%[ma:9e233b98]
 Trên một giá sách  có $15$ quyển sách, trong đó có $5$ quyển văn nghệ. Lấy ngẫu nhiên từ đó ba quyển. Khi đó
  \choiceTF
  {\True Xác suất để lấy ngẫu nhiên $3$ quyển trong đó có $1$ cuốn  văn nghệ là $\dfrac{45}{91}$}
  {Xác suất để lấy ngẫu nhiên $3$ quyển trong đó  có $2$ cuốn văn nghệ là $\dfrac{14}{91}$}
  {\True Xác suất để lấy ngẫu nhiên $3$ quyển trong đó có $3$ cuốn văn nghệ  là $\dfrac{2}{9}$}
  {\True Xác suất sao cho có ít nhất một quyển  văn nghệ là $\dfrac{67}{91}$}
  \loigiai{
  Số phần  tử của không gian mẫu $n(\Omega)=\mathrm{C}_{15}^3$.
  \begin{itemchoice}
    \itemch Gọi biến cố $A\colon$ \lq\lq Lấy $3$ quyển trong đó có $1$ cuốn văn nghệ\rq\rq.\\
    Suy ra $n(A)=\mathrm{C}_5^1\cdot \mathrm{C}_{10}^2$.\\
   Xác suất để lấy ngẫu nhiên  $3$ quyển trong đó có $1$ cuốn văn nghệ là $$\mathrm{P}(A)=\dfrac{\mathrm{C}_5^1\cdot \mathrm{C}_{10}^2}{\mathrm{C}_{15}^3}=\dfrac{45}{91}.$$
   \itemch Gọi biến cố $B\colon$ \lq\lq Lấy  $3$ quyển trong đó có $2$ cuốn văn nghệ\rq\rq.\\
   Suy  ra $n(B)=\mathrm{C}_5^2\cdot \mathrm{C}_{10}^1$.\\
   Xác suất để lấy ngẫu nhiên $3$ quyển trong đó có $2$ cuốn văn nghệ là $$\mathrm{P}(B)=\dfrac{\mathrm{C}_5^2\cdot  \mathrm{C}_{10}^1}{\mathrm{C}_{15}^3}=\dfrac{20}{91}.$$
   \itemch Gọi biến cố $C\colon$ \lq\lq  Lấy $3$ quyển trong đó có $3$ cuốn văn nghệ\rq\rq.\\
    Suy ra $n(C)=\mathrm{C}_5^3$.\\
   Xác suất để  lấy ngẫu nhiên $3$ quyển trong đó có $3$ cuốn văn nghệ là $$\mathrm{P}(C)=\dfrac{\mathrm{C}_5^3}{\mathrm{C}_{15}^3}=\dfrac{2}{91}.$$
   \itemch Xác suất để lấy ngẫu nhiên $3$ quyển trong đó có ít nhất $1$ cuốn văn  nghệ là
   $$\mathrm{P}=\mathrm{P}(A)+\mathrm{P}(B)+\mathrm{P}(C)=\dfrac{2}{91}+\dfrac{45}{91}+\dfrac{20}{91}=\dfrac{67}{91}.$$ 
  \end{itemchoice}
 } 
\end{ex}

% nguồn: Ok_dungsai_lop11_Lan1
\begin{ex}%[1D9V2-5]%[ma:48662926]
 Cả hai xạ thủ cùng bắn vào bia. Xác suất người thứ nhất bắn trúng bia là $0{,}8$; người thứ hai bắn trúng bia là $0{,}7$. Khi đó  xác suất để:
  \choiceTF
  {Người thứ nhất bắn trúng và người thứ hai bắng không trúng bia bằng  $0{,}14$}
  {\True Người thứ nhất bắn không trúng  và người thứ hai bắn trúng bia bằng $0{,}14$}
  {\True Hai người đều bắn trúng bia bằng  $0{,}56$}
  {\True  Có ít nhất một người bắn trúng bia bằng $0{,}94$}
  \loigiai{
  Gọi $A$ là biến cố \lq\lq  Người thứ nhất bắn trúng bia\rq\rq. Suy ra $\mathrm{P}(A)=0{,}8$, $\mathrm{P}(\overline{A})=0{,}2$.\\
   Gọi $B$ là biến cố \lq\lq Người thứ hai bắn trúng bia\rq\rq. Suy ra $\mathrm{P}(B)=0{,}7$, $\mathrm{P}(\overline{B})=0{,}3$.
  \begin{itemchoice}
   \itemch Biến cố  người thứ nhất bắn trúng và người thứ hai bắng không trúng bia là $A\overline{B}$ và $\mathrm{P}(A\overline{B})=\mathrm{P}(A)\cdot \mathrm{P}(\overline{B})=0{,}8\cdot 0{,}3=0{,}24$.
   \itemch Biến cố người thứ nhất bắn không trúng và người thứ hai bắn trúng bia là $\overline{A}B$ và $\mathrm{P}(\overline{A}B)=\mathrm{P}(\overline{A})\cdot \mathrm{P}(B)=0{,}2\cdot  0{,}7=0{,}14$.
   \itemch  Biến cố cả hai người đều bắn trúng bia là $A B$ và $\mathrm{P}(A B)=\mathrm{P}(A)\cdot \mathrm{P}(B)=0{,}8\cdot 0{,}7=0{,}56$.
   \itemch Gọi $C$ là biến cố \lq\lq Có ít nhất một người bắn trúng bia\rq\rq. Xác suất để có ít  nhất một người bắn trúng là\\
   $\mathrm{P}(C)=\mathrm{P}(\overline{A B})+\mathrm{P}(\overline{A}B)+\mathrm{P}(A  B)=0{,}24+0{,}14+0{,}56=0{,}94$.
  \end{itemchoice}
  }
\end{ex}


%%%%%%%%%%%%%%%% TRẢ LỜI NGẮN %%%%%%%%%%%%%%%%

% nguồn: Ok_TraLoiNgan_Lop11_Lan1
\begin{ex}%[1D9N2-5]%[ma:ba2056e1]
	 Chọn ngẫu nhiên $3$ đỉnh trong số $12$ đỉnh của một đa giác đều $12$ cạnh. Tính xác suất của biến cố $A$ \lq\lq $3$ đỉnh được chọn tạo thành một tam giác vuông\rq\rq(làm tròn đến hàng phần trăm).
	 \shortans{0,06}
	 \loigiai{
		Số tam giác có thể được tạo ra bằng cách lấy $3$ đỉnh trong $12$ đỉnh  của đa giác là $\mathrm{C}_{12}^3$.\\
		Vì đa giác là đa giác đều nên sẽ  có $1$ đường tròn ngoại tiếp đa giác đó.\\
		Để $3$  đỉnh lấy ra tạo thành một tam giác vuông thì phải có $2$ đỉnh tạo thành đường kính của đường tròn ngoại tiếp. Đa giác đó số đường kính là $6$.\\
		Số cách chọn ra $3$ đỉnh từ  $12$ đỉnh của đa giác để tạo thành tam giác vuông là $n(A)=\mathrm{C}_{10}^1\cdot6$.\\
	Xác suất của biến cố  $A$ là $\mathrm{P}(A)=\dfrac{\mathrm{C}_{10}^1\cdot6}{\mathrm{C}_{12}^3}=\dfrac{5}{77}$.}
\end{ex}

% nguồn: Ok_TraLoiNgan_Lop11_Lan1
\begin{ex}%[1D9H2-5]%[ma:b278ec87]
	Gieo ngẫu nhiên $3$ con xúc xắc cân đối và đồng chất. Tính xác suất của biên cố $A\colon$ `` Tích số chấm xuất hiện trên mỗi con  xúc xắc chia hết cho $15$''.
	 \shortans{0,25}
	 \loigiai{
		Gọi biến cố $B\colon $`` Tích số chấm xuất hiện trên mỗi con xúc sắc không chia hết cho $5$'', biến cố $C\colon $`` Tích số chấm xuất  hiện trên mỗi con xúc sắc không chia hết cho $3$''.\\
		Khi đó $A$ là biến cố đối của biến cố $B\cup  C$.\\
		Ta có $P(B)=\dfrac{5}{6}  \cdot \dfrac{5}{6}\cdot \dfrac{5}{6}=\left(\dfrac{5}{6}\right)^3$, $P(C)=\dfrac{4}{6} \cdot \dfrac{4}{6}\cdot \dfrac{4}{6} =\left(\dfrac{4}{6}\right)^3$, $P(BC)=\dfrac{3}{6}\cdot \dfrac{3}{6}\cdot \dfrac{3}{6}=\left(\dfrac{3}{6}\right)^3$. Khi đó
		\[P(B\cup C)=P(B)+P(C)-P(BC)=\left(\dfrac{5}{6}\right)^3+\left(\dfrac{5}{6}\right)^3-\left(\dfrac{3}{6}\right)^3=\dfrac{3}{4}.\] 
		Do đó xác suất của biến cố $A$ là  $P(A)=1-P(B\cup C)=1-\dfrac{3}{4}=\dfrac{1}{4}$.
	} 
\end{ex}

% nguồn: Ok_TraLoiNgan_Lop11_Lan1
\begin{ex}%[1D9H2-5]%[ma:17e24c34]
	Tại một trường trung học phổ thông $X$, có $12\%$ học sinh học giỏi môn Tiếng Anh, $35\%$ học sinh  học giỏi môn Toán và $8\%$ học sinh học giỏi cả hai môn Toán, Tiếng Anh. Chọn ngẫu nhiên một học sinh từ trường $X$, tính xác suất để chọn được một học sinh không giỏi môn nào trong hai môn Toán, Tiếng Anh.
	\shortans{0,61} 
	 \loigiai{
		Gọi $A$ là biến cố \lq\lq Chọn được một học sinh giỏi môn  Tiếng Anh\rq\rq, $B$ là biến cố \lq\lq Chọn được một học sinh giỏi môn Toán\rq\rq.\\
		Xác suất để chọn được  một học sinh giỏi Toán hoặc giỏi Anh là
		\begin{eqnarray*}
			\mathrm{P}(A\cup B)&=&  \mathrm{P}(A)+\mathrm{P}(B)-\mathrm{P}(AB)\\
			&=&\dfrac{12}{100}+\dfrac{35}{100}-\dfrac{8}{100}\\ 
			&=&0{,}39. 
		\end{eqnarray*}
		Xác suất để chọn được một em học sinh không giỏi môn nào trong  hai môn Toán, Tiếng Anh là
	$$\mathrm{P}(\overline{A\cup B})=1-\mathrm{P}(A\cup B)=1-0{,}39=0{,}61.$$} 
\end{ex}

% nguồn: Ok_TraLoiNgan_Lop11_Lan1
\begin{ex}%[1D9H2-5]%[ma:c11ea68d]
	Người ta thăm dò một số lượng người hâm mộ bóng đá tại một thành phố, nơi có hai đội bóng đá $X$ và $Y$ cùng thi đấu giải vô địch quốc gia. Biết rằng số lượng người hâm mộ đội bóng đá $X$ là $22\%$, số lượng người hâm mộ đội bóng đá $Y$ là $39\%$, trong số đó có $7\%$ người nói rằng họ  hâm mộ cả hai đội bóng trên. Chọn ngẫu nhiên một người hâm mộ trong số những người được hỏi, tính xác suất để chọn được người không hâm mộ đội nào trong hai đội bóng đá $X$ và $Y$.
	\shortans{0,46} 
	\loigiai{ 
		Gọi $A$ là biến cố \lq\lq Chọn được một người hâm mộ đội  bóng đá $X$\rq\rq, $B$ là biến cố \lq\lq Chọn được một người hâm mộ đội bóng đá $Y$\rq\rq.\\
		 Khi đó
		\begin{itemize}
			\item  $\mathrm{P}(A)=\dfrac{22}{100}=0{,}22$;
			\item  $\mathrm{P}(B)=\dfrac{39}{100}=0{,}39$;
			\item $ \mathrm{P}(A B)=\dfrac{7}{100}=0{,}07$. 
		\end{itemize}
		Suy ra $\mathrm{P}(A\cup B)=\mathrm{P}(A)+\mathrm{P}(B)-\mathrm{P}(A  B)=0{,}22+0{,}39-0{,}07=0{,}54$.\\
		Xác suất để chọn được người không hâm mộ đội nào trong hai đội bóng đá $X$ và  $Y$ là
	$$\mathrm{P}(\overline{A\cup B})=1-\mathrm{P}(A\cup B)=1-0{,}54=0{,}46.$$} 
\end{ex}

% nguồn: Ok_TraLoiNgan_Lop11_Lan1
\begin{ex}%[1D9H2-5]%[ma:7c19490d]
	Rút ngẫu nhiên $1$ lá bài từ bộ bài tây $52$ lá. Tính xác suất của biến cố \lq\lq Lá bài được chọn có màu đen hoặc lá đó có số  chia hết cho $3$\rq\rq(làm tròn đến hàng phần trăm).
	\shortans{0,62} 
	 \loigiai{
		Số phần tử của không gian mẫu  $n(\Omega)=52$.\\
		Gọi $A$ là biến cố \lq\lq Lá bài được chọn có màu đen\rq\rq\,suy  ra $n(A)=26$ hay $\mathrm{P}(A)=\dfrac{26}{52}=\dfrac{1}{2}$.\\
		$B$ biến cố \lq\lq Lá bài được chọn  có số chia hết cho $3$\rq\rq\,suy ra $n(B)=12$ hay $\mathrm{P}(A)=\dfrac{12}{52}=\dfrac{3}{13}$.\\
		Vì  $A$ và $B$ là hai biến cố độc lập nên ta có
	$$\mathrm{P}(A\cup B)=\mathrm{P}(A)+\mathrm{P}(B)-\mathrm{P}(A B)=\dfrac{1}{2}+\dfrac{1}{13}-\dfrac{1}{2}\cdot\dfrac{1}{13}=\dfrac{8}{13}.$$} 
\end{ex}

% nguồn: Ok_TraLoiNgan_Lop11_Lan1
\begin{ex}%[1D9V2-5]%[ma:fc494f99]
	Hai bạn An và Bình cùng tập ném bóng rổ một cách độc lập ở hai nửa sân khác nhau. Xác suất bạn An và bạn Bình ném bóng vào rổ lần lượt là $0{,}6$ và $0{,}9$. Trong cùng một lần  ném, tính xác suất có ít nhất một bạn ném bóng vào rổ.
	\shortans{0,96} 
	\loigiai{ 
		Xét các biến cố  $A$: \lq\lq Bạn An ném bóng trúng rổ\rq\rq;
		$B$: \lq\lq  Bạn Bình ném bóng trúng rổ\rq\rq;
		$C$: \lq\lq  Có ít nhất một bạn ném bóng vào rổ\rq\rq.\\
		$A$ và $B$ là hai biến cố độc lập nên $\mathrm{P}\left(A\cap  B\right)=\mathrm{P}\left(A\right)\cdot \mathrm{P}\left(B\right)$.\\
	Khi đó $\mathrm{P}\left(C\right)=\mathrm{P}\left(A\cup B\right)=\mathrm{P}\left(A\right)+\mathrm{P}\left(B\right)-\mathrm{P}\left(B\right)\cdot \mathrm{P}\left(C\right)=0{,}6+0{,}9-0{,}6\cdot  0{,}9=0{,}96$.}
\end{ex}

% nguồn: Ok_TraLoiNgan_Lop11_Lan1
\begin{ex}%[1D9C2-5]%[ma:6580955c]
	Một hộp chứa 3 quả bóng xanh và một số quả bóng đỏ có cùng kích thước và khối lượng. Lấy ra ngẫu nhiên 2 quả bóng từ hộp. Biết rằng xác suất của biến cố \text{\lq\lq}Lấy được 2 quả bóng đỏ\text{\rq\rq} gấp 5 lần xác suất của biến cố \text{\lq\lq}Lấy được 2 quả bóng xanh\text{\rq\rq}. Tính  xác suất của biến cố \text{\lq\lq}Lấy được 2 quả bóng có cùng màu\text{\rq\rq}.
	 \shortans{0,5}
	 \loigiai{
	Gọi $n$ là số bóng đỏ có trong hộp. Khi đó, tổng số bóng có trong hộp là $n+3$. Xác suất lấy được 2  quả bóng xanh là: $\dfrac{\mathrm{C}_3^2}{\mathrm{C}_{n+3}^2}$.\\
	Xác  suất lấy được 2 quả bóng đỏ là: $\dfrac{\mathrm{C}_n^2}{\mathrm{C}_{n+3}^2}$. Theo đề bài, ta có:
	 \[\dfrac{\mathrm{C}_n^2}{ \mathrm{C}_{n+3}^2}=5 \cdot\dfrac{ \mathrm{C}_3^2}{ \mathrm{C}_{n+3}^2} \Leftrightarrow\dfrac{n(n-1)}{2}=15.\]
	Do $n$ là  số tự nhiên nên $n=6$.\\
	Do  đó, xác suất của biến cố \text{\lq\lq}Cả 2 quả bóng lấy ra đều có cùng màu\text{\rq\rq} là
	\[\dfrac{  \mathrm{C}_3^2}{ \mathrm{C}_{n+3}^2}+\frac{\mathrm{C}_n^2}{ \mathrm{C}_{n+3}^2}=\dfrac{ \mathrm{C}_3^2}{ \mathrm{C}_9^2}+\dfrac{ \mathrm{C}_6^2}{ \mathrm{C}_9^2}=\dfrac{1}{2}=0{,}5.\]
	 }
\end{ex}
